%! TEX root = ../m1-19-20.tex

\chapter{Examen 2018--2019}

\section{Numărul 1}

\indent\indent 1. Să se afle valorile extreme ale funcției:
\[
  f(x,y) = x^2 + y^2 - 3x - 2y + 1
\]
pe mulțimea $K: x^2 + y^2 \leq 1$.

\vspace{0.5cm}

2. Calculați $\ds\int_0^\infty \frac{dx}{x^4 + 1}$.

\vspace{0.5cm}

3. Să se calculeze volumul corpului mărginit de suprafețele:
\[
  \begin{cases}
    z &= x^2 + y^2 \\
    2 &= x^2 + y^2 + z^2.
  \end{cases}
\]

\vspace{0.5cm}

4. Să se calculeze fluxul cîmpului $\vec{V} = (x + z^2, y + z^2, -2z)$
prin suprafața $x^2 + y^2 + z^2 = 1, z \geq 0$.

\vspace{0.5cm}

5. Să se calculeze aria suprafeței:
\[
  \Sigma = \{ (x, y, z) \mid x^2 + y^2 + z^2 = R^2, z \geq 0 \} \cap %
  \{ (x, y, z) \mid x^2 + y^2 \leq Ry \}.
\]
(Aria lui Viviani)

%%%%%%%%%%%%%%%%%%%%%%%%%%%%%%%%%%%%%%%%%%%%%%%%%%%%%%%%%%%%%%%%%%%%%%

\section{Numărul 2}

\indent\indent 1. Să se afle valorile extreme ale funcției:
\[
  f(x, y) = x \cdot y, \text{ pe } K: 2x^2 + y^2 \leq 1.
\]

\vspace{0.5cm}

2. Calculați $\ds\int_0^\infty \frac{dx}{x^3 + 1}$.

\vspace{0.5cm}

3. Să se calculeze volumul corpului obținut prin intersecția suprafețelor:
\[
  \begin{cases}
    2 &= x^2 + y^2 + z^2 \\
    z^2 &= x^2 + y^2, z \geq 0
  \end{cases}
\]

\vspace{0.5cm}

4. Să se calculeze aria suprafeței $2 - z = x^2 + y^2, z \in [0, 1]$.

\vspace{0.5cm}

5. Calculați fluxul cîmpului vectorial $\vec{V} = \frac{q}{4 \pi} \cdot \frac{\vec{r}}{r^3}$
prin:
\begin{enumerate}[(a)]
\item Sfera centrată în origine și cu rază $R > 0$;
\item Orice suprafață $\Sigma$ închisă, care nu conține originea.
\end{enumerate}
Observație: $\vec{r} = (x, y, z)$ și $r = ||\vec{r}|| = \sqrt{x^2 + y^2 + z^2}$.

(Legea lui Gauss)

%%%%%%%%%%%%%%%%%%%%%%%%%%%%%%%%%%%%%%%%%%%%%%%%%%%%%%%%%%%%%%%%%%%%%%
\section{Restanță}


\indent\indent 1. Studiați convergența seriei de puteri:
\[
  \sum_{n \geq 1} \sqrt{n}{4^n + 3} x^n.
\]

\vspace{0.5cm}

2. Studiați convergența uniformă a șirului de funcții:
\[
  f_n : [3, \infty) \to \RR, \quad f_n(x) = \frac{n}{x(x^2 + n)}, n \geq 1.
\]

\vspace{0.5cm}

3. Calculați volumul mulțimii mărginite de suprafețele:
\begin{align*}
  S_1 &: x^2 + y^2 + z^2 = 6z, z \geq 3 \\
  S_2 &: 2(x^2 + y^2) = z^2.
\end{align*}

\vspace{0.5cm}

4. Fie suprafața $ S = \{ x^2 + y^2 + z^2 = 5, z \leq 0 \} $.
Desenați suprafața și calculați:
\[
  \iint_S (2x + 5x) dy \wedge dz + 2ydz \wedge dx + (2z - 5x) dx \wedge dy.
\]

